\documentclass[11pt]{article}
\usepackage[utf8]{inputenc}
\usepackage[T1]{fontenc}
\usepackage[margin=0.82in,headheight=15pt]{geometry}
\usepackage{lmodern}
\usepackage{amsmath,amssymb}
\usepackage{array}
\usepackage{booktabs}
\usepackage{longtable}
\usepackage{enumitem}
\usepackage{xcolor}
\usepackage[hidelinks]{hyperref}
\usepackage{titlesec}
\usepackage{fancyhdr}

\setlength{\parindent}{0pt}
\setlength{\parskip}{0.42em}
\setlist[itemize]{nosep,leftmargin=1.35em}
\setlist[enumerate]{nosep,leftmargin=1.55em}
\definecolor{Accent}{HTML}{0B4F6C}
\definecolor{Soft}{HTML}{F2F6FA}
\titleformat{\section}{\large\bfseries\color{Accent}}{}{0pt}{}
\titleformat{\subsection}{\normalsize\bfseries\color{Accent}}{}{0pt}{}
\pagestyle{fancy}
\fancyhf{}
\fancyhead[L]{Chaos Toolbox}
\fancyhead[R]{Explorador Sprott}
\fancyfoot[C]{\thepage}
\newcommand{\code}[1]{\texttt{#1}}
\newcommand{\R}{\mathbb{R}}

\begin{document}

{\LARGE\bfseries Explorador Sprott}\par
{\small Guia matematica y de uso para la reimplementacion educativa inspirada en Julien C. Sprott.}\par
\vspace{0.8em}
\tableofcontents
\newpage

\section{Que hace esta pestana}

El Explorador Sprott no es una coleccion fija de sistemas como Lorenz o Rossler. Es un
laboratorio de \emph{ecuaciones compactas}: un codigo corto representa una familia de
mapas o flujos polinomiales, la toolbox lo decodifica, lo simula con el backend C y
clasifica rapidamente el resultado.

El flujo correcto es:
\begin{enumerate}
  \item elegir un codigo de ejemplo o cargar un diccionario local \code{.DIC};
  \item decodificarlo para ver que familia, dimension, orden y coeficientes representa;
  \item simularlo y descartar un transitorio;
  \item mirar la proyeccion graficada;
  \item interpretar la etiqueta: divergente, punto fijo, baja complejidad o candidato.
\end{enumerate}

La etiqueta \emph{candidate\_chaotic} significa solamente que la trayectoria quedo acotada
y no colapso en los filtros rapidos. No es una prueba formal de caos.

\section{Decodificador}

Un codigo tipo Sprott se separa en dos partes:
\[
\boxed{\text{letra de familia}}\quad\boxed{\text{caracteres de coeficientes}}
\]

La primera letra define el tipo de ecuacion:
\begin{longtable}{p{0.16\linewidth}p{0.22\linewidth}p{0.18\linewidth}p{0.26\linewidth}}
\toprule
Letras & Tipo & Dimension & Ordenes \\
\midrule
\endhead
A--D & mapa polinomial & 1 & 2, 3, 4, 5 \\
E--H & mapa polinomial & 2 & 2, 3, 4, 5 \\
I--L & mapa polinomial & 3 & 2, 3, 4, 5 \\
M--P & mapa polinomial & 4 & 2, 3, 4, 5 \\
Q--T & flujo polinomial & 3 & 2, 3, 4, 5 \\
U--X & flujo polinomial & 4 & 2, 3, 4, 5 \\
Y--Z & funciones especiales & pendiente & pendiente \\
\bottomrule
\end{longtable}

Los coeficientes se decodifican inicialmente con
\[
c=\frac{\operatorname{ord}(\chi)-77}{10}.
\]
Asi, \(\chi=M\) produce \(c=0\), letras antes de \(M\) producen coeficientes negativos
y letras despues de \(M\) producen coeficientes positivos. Los diccionarios historicos
tambien pueden contener otros caracteres imprimibles; esta reimplementacion los acepta
solo como datos locales cargados por el usuario.

\section{Mapas}

Un mapa discreto actualiza el estado por iteraciones:
\[
x_{n+1}=F(x_n), \qquad x_n\in\R^D.
\]
En una familia polinomial, cada componente usa la misma base de monomios:
\[
F_i(x)=\sum_{j=1}^{N_m} c_{ij}m_j(x).
\]
El numero de monomios de grado total menor o igual a \(O\) en dimension \(D\) es
\[
N_m=\binom{D+O}{O},
\]
y el numero total de coeficientes es
\[
N_c=D\binom{D+O}{O}.
\]

Para ver figuras interesantes en mapas:
\begin{itemize}
  \item usa ejemplos de \code{SELECTED.DIC} o \code{BOOKFIGS.DIC};
  \item usa al menos 8\,000 a 20\,000 iteraciones;
  \item descarta 1\,000 a 3\,000 puntos de transitorio;
  \item mira primero la proyeccion \(x\)-\(y\).
\end{itemize}

\section{Flujos}

Un flujo continuo define una ecuacion diferencial:
\[
\dot{x}=F(x).
\]
La toolbox puede integrar estos sistemas con Euler o RK4. Euler es util para reproducir
el espiritu de programas historicos; RK4 suele ser mas estable para exploracion moderna:
\[
x(t+h)\approx x(t)+hF(x(t)) \quad \text{(Euler)}.
\]

Para flujos, reduce \(h\) si la trayectoria explota. Si no aparece estructura, aumenta
el numero de pasos antes de concluir.

\section{Exploracion y botones}

\begin{longtable}{p{0.24\linewidth}p{0.64\linewidth}}
\toprule
Control & Uso recomendado \\
\midrule
\endhead
Preset & Carga parametros razonables. Empieza por mapas 2D rapidos. \\
Tipo & \code{map} itera directamente; \code{flow} integra una EDO. \\
Dimension & Numero de variables. En la grafica se muestra \(x\)-\(y\) si hay dos o mas. \\
Orden & Grado maximo del polinomio. Orden alto aumenta mucho el numero de coeficientes. \\
Iteraciones/pasos & Mas puntos pueden revelar estructura, pero tardan mas. \\
Transitorio & Puntos descartados al inicio. Evita graficar el arranque. \\
Tolerancia divergencia & Si la norma supera ese umbral, se detiene y se marca divergente. \\
Generar codigo & Crea un codigo sintetico. Es didactico, no garantiza atractor interesante. \\
Simular y graficar & Ejecuta el codigo actual y actualiza la grafica. \\
Buscar candidato & Si hay \code{.DIC} local cargado, prueba codigos del libro primero; si no, usa aleatorios sinteticos. \\
\bottomrule
\end{longtable}

\section{Ejemplos locales del libro}

Los archivos \code{BOOKFIGS.DIC}, \code{SELECTED.DIC} y \code{SPECIAL.DIC} no se
redistribuyen en el repositorio. La pestana puede leerlos desde una carpeta local que
el usuario ya tenga descargada.

Uso recomendado:
\begin{enumerate}
  \item abrir \emph{Ejemplos};
  \item revisar la columna \emph{Ejemplos del libro desde .DIC local};
  \item si no aparece nada, pulsar \emph{Elegir .DIC};
  \item seleccionar \code{SELECTED.DIC} o \code{BOOKFIGS.DIC};
  \item pulsar \emph{Cargar codigos};
  \item seleccionar una linea y pulsar \emph{Simular codigo local}.
\end{enumerate}

\section{Atribucion y limites legales}

Inspirado por Julien C. Sprott, \emph{Strange Attractors: Creating Patterns in Chaos},
M\&T Books, 1993.

Esta toolbox es una reimplementacion educativa independiente, no una version oficial.
El software original, diccionarios, ejecutables, figuras y archivos del disco permanecen
copyright de J. C. Sprott y no se redistribuyen en este repositorio.

\end{document}
